\documentclass[11pt]{article}
\usepackage[margin=1in]{geometry}
\usepackage{booktabs}
\usepackage{enumitem}
\usepackage{hyperref}
\usepackage{xcolor}
\usepackage{mdframed}
\usepackage{array}

\definecolor{concern}{RGB}{180,0,0}
\definecolor{ok}{RGB}{0,128,0}
\definecolor{question}{RGB}{0,80,160}

\newcommand{\concern}[1]{\textcolor{concern}{\textbf{Concern:} #1}}
\newcommand{\ok}[1]{\textcolor{ok}{\textbf{OK:} #1}}
\newcommand{\question}[1]{\textcolor{question}{\textbf{Open question:} #1}}

\title{\textbf{Deep Review: Editor's Letter}\\
\large The Rise of Reddit --- Management Science Resubmission}
\date{Reviewed September 2026}

\begin{document}
\maketitle

This document records a close reading of the editor's three comments and
the authors' responses, cross-checked against the revised text and tables.

\tableofcontents
\newpage

%----------------------------------------------------------------------
\section{Editor Comment 1: Use Pedersen (2022) as Conceptual Motivation}
%----------------------------------------------------------------------

\subsection*{The Editor's Request}
\begin{mdframed}
\textit{``The AE communicated to me that unless you are prepared to engage
more deeply with the model, it may be more effective to use Pedersen (2022)
as a conceptual motivation for empirical tests, rather than positioning the
paper as a direct test of the model. I completely agree. I think you should
emphasize more the empirical aspects of your study.''}
\end{mdframed}

\subsection*{What the Authors Claim to Have Done}
Revised Abstract, Section~1, Section~2 (renamed ``Conceptual Framework and
Empirical Implications''), and Section~7 to downplay direct-test language
and use ``motivated by,'' ``inspired by,'' and ``drawing on'' instead.

\subsection*{What Is Actually in the Text}

\paragraph{What changed.}
\begin{itemize}
  \item Section~2 is renamed ``Conceptual Framework and Empirical
    Implications'' --- a clear and appropriate signal.
  \item The opening of Section~2 now reads: ``We briefly summarize the
    features of the model most relevant to our setting and draw on its
    underlying economic mechanisms to \textbf{motivate and organize} our
    empirical analysis.''
  \item The contributions paragraph in Section~1 uses ``inspired by the
    guiding framework'' and ``motivated by Pedersen (2022)'' rather than
    ``we test Proposition~X.''
  \item The abstract opens with ``Motivated by Pedersen (2022)\ldots''
\end{itemize}

\paragraph{What has not changed.}
The six empirical hypotheses (H1--H6) remain tightly pinned to specific
propositions of Pedersen (2022):

\begin{itemize}
  \item H1 and H2: ``Motivated by this mechanism [Proposition~1]\ldots''
  \item H3 and H4: ``Motivated by Propositions~4 and~5\ldots''
  \item H5 and H6: ``motivated by Propositions~7 and~8\ldots''
\end{itemize}

\subsection*{Assessment}

\concern{The hypothesis framing is still tightly coupled to specific
propositions.} Calling something ``motivated by Proposition~4'' is
functionally close to saying ``this is a test of Proposition~4.'' An
editor or AE reading carefully could feel the paper has changed its
language but not its posture. The concern was not just about word choice
--- it was about whether the empirical design can actually deliver a
convincing test of the model's mechanisms. The language revision helps at
the margin, but the one-to-one mapping between H1--H6 and Pedersen's
propositions leaves the paper exposed to the same critique.

\question{Is the hypothesis structure tight enough to the propositions that
it could invite the same pushback in the next round? Consider whether
hypotheses can be motivated more generically (e.g., from broad social
network intuition) with Pedersen cited as one formalization, rather than
deriving each hypothesis from a named proposition.}

\newpage
%----------------------------------------------------------------------
\section{Editor Comment 2: Differentiate Types of Social Media Interactions}
%----------------------------------------------------------------------

\subsection*{The Editor's Request}
\begin{mdframed}
\textit{``Would it be possible to differentiate the type of social media
interactions? It makes a huge difference seeing an anonymous post versus
seeing Trump's tweets. You would need to uncover deeper insights and more
interesting patterns from the data to provide sufficient contribution to
the literature.''}
\end{mdframed}

\subsection*{What the Authors Claim to Have Done}
Introduced a celebrity vs.\ other user split in Section~5.4 and Table~9,
defining celebrity users as those whose average PageRank exceeds the weekly
90th~percentile.

\subsection*{What Is Actually in the Text and Tables}

\paragraph{What works.}
\begin{itemize}
  \item Celebrity user tone has stronger effects on returns: rational 2~bps,
    fanatic 6~bps, na\"{i}ve 3~bps vs.\ 1~bp, 2~bps, 2~bps for other users.
  \item Celebrity fanatic tone has a particularly large retail flow effect
    (0.66\%) vs.\ other fanatic users (0.24\%) --- an interesting and
    economically meaningful difference.
  \item The economic narrative is coherent: network prominence amplifies
    the salience of fanatic views and triggers retail trading.
\end{itemize}

\paragraph{What is missing from the discussion.}
Table~9 shows that \textbf{short-flow effects are entirely null for both
celebrity and other users} --- all Granger causality $p$-values for
ShortFlow(t) exceed 57\% (celebrity: 68.6\%, 63.6\%, 57.8\%; other:
98.7\%, 74.3\%, 5.4\%). The text does not mention this at all.

This is notable because the main results (Tables~5 and~7) show a strong,
significant negative relation between tone and short flow in the
high-influence subsample. Yet once the sample is split by celebrity status
rather than aggregate influence, the short-flow relation disappears
entirely. This deserves at least a sentence of explanation.

\concern{The null short-flow result for celebrity users is unreported and
potentially inconsistent with the main findings. Leaving it unacknowledged
may invite reviewer questions in the next round.}

\question{Is there an economic story for why short sellers do not
differentially respond to celebrity vs.\ non-celebrity tone? One
possibility: the celebrity split uses the \textit{whole} sample (not the
high-influence subsample), so the average influence level is lower, which
could explain the null. If so, a brief note to this effect would preempt
the question.}

\newpage
%----------------------------------------------------------------------
\section{Editor Comment 3: Technological Progress and Competing Platforms}
%----------------------------------------------------------------------

\subsection*{The Editor's Request}
\begin{mdframed}
\textit{``Can you relate the discussion to technological progress? Would
the emergence of competing platforms affect the results? I am just trying
to suggest something that would lend a good angle for contribution to the
literature, in addition to the prediction results.''}
\end{mdframed}

\subsection*{What the Authors Claim to Have Done}
\begin{itemize}
  \item Section~6.1 (Table~10): ChatGPT release (November~30, 2022) as
    technological cutoff; PVAR estimated separately pre- and post-ChatGPT.
  \item Section~6.2 (Table~11): Relative activity on StockTwits, X, and
    Seeking Alpha used to define platform-emergence periods (2020--2023).
\end{itemize}

\subsection*{What Is Actually in the Text and Tables}

\subsubsection*{3.1 Technological Progress (Section 6.1 / Table 10)}

\paragraph{What works.}
The core finding is clean and interesting: agent tones predict returns
only in the pre-ChatGPT period (Table~10 Panel~A). For example, the
RationalTone coefficient on Return(t) falls from $0.0014^{***}$
($t=4.48$) pre-ChatGPT to $0.0009$ ($t=1.02$, insignificant)
post-ChatGPT. The explanation --- AI lowers content-production costs,
diluting the information content of posts --- is economically intuitive
and well-grounded in cited evidence.

\paragraph{A striking result worth highlighting more.}
Table~10 Panel~B shows a \textbf{sign flip in the short-flow relation}
around ChatGPT:

\begin{center}
\begin{tabular}{lcc}
\toprule
 & Pre-ChatGPT & Post-ChatGPT \\
\midrule
RationalTone $\to$ ShortFlow & $-0.0260^{***}$ [$-3.17$] & $+0.0190^{***}$ [$+2.80$] \\
FanaticTone $\to$ ShortFlow  & $-0.0343^{***}$ [$-2.58$] & $+0.0216^{***}$ [$+2.60$] \\
Na\"{i}veTone $\to$ ShortFlow & $-0.0258^{***}$ [$-5.19$] & $+0.0088^{**}$ [$+2.52$] \\
\bottomrule
\end{tabular}
\end{center}

Pre-ChatGPT, short sellers back off from bullish tone (consistent with
main results). Post-ChatGPT, short sellers \textit{increase} shorting in
response to bullish tone --- i.e., they now trade \textit{against} social
media sentiment rather than accommodating it. All six coefficients are
statistically significant. The text offers an explanation (``as the
informational value of social-media content for asset prices weakens
after ChatGPT, more sophisticated investors increasingly trade against
the sentiment''), which is plausible and internally consistent with the
null return predictability post-ChatGPT.

\ok{The sign flip is robust and the explanation is coherent. This is one
of the more striking results in the paper and is currently underemphasized.
Consider elevating it in the abstract or introduction as a novel finding.}

\subsubsection*{3.2 Competing Platforms (Section 6.2 / Table 11)}

\paragraph{What works.}
Belief transmission weakens during platform-emergence periods: rational
and fanatic tone become insignificant predictors of na\"{i}ve tone
(Table~11 Panel~A, $t=1.47$ and $t=1.73$ respectively). Return
predictability also disappears (all insignificant). During other periods
(Panel~B), all effects are strongly significant. The contrast is clean.

\paragraph{Potential inconsistency with the text.}
Section~6.2 states that ``this attenuation also extends to subsequent
return and \textbf{trading dynamics}.'' However, in Table~11 Panel~A
(emergence periods):

\begin{center}
\begin{tabular}{lcc}
\toprule
 & Emergence periods & Other periods \\
\midrule
RationalTone $\to$ RetailFlow & $0.0115^{**}$ [$2.16$], $p=3.1\%$ & $0.0176^{***}$ [$6.93$] \\
FanaticTone $\to$ RetailFlow  & $0.0008$ [$0.15$], $p=87.7\%$ & $0.0228^{***}$ [$6.01$] \\
Na\"{i}veTone $\to$ RetailFlow & $0.0077^{**}$ [$2.10$], $p=3.6\%$ & $0.0098^{***}$ [$5.98$] \\
\bottomrule
\end{tabular}
\end{center}

Fanatic tone's effect on retail flow does weaken (as the text implies),
but rational and na\"{i}ve tone still significantly predict retail flow
at the 5\% level during emergence periods. The claim of broad attenuation
in trading dynamics is therefore only partially supported.

\concern{The text overstates the attenuation in retail trading during
platform-emergence periods. Rational and na\"{i}ve tone effects on retail
flow remain significant. This should be qualified --- e.g., ``the
weakening is most pronounced for fanatic tone, while rational and na\"{i}ve
tone retain some predictive power for retail trading.'' Without this
qualification, a careful reviewer may flag the discrepancy between text
and table.}

\end{document}
